Rainforest climate monitoring loses observation value when onboard controllers trigger captures on a fixed nadir schedule while clouds pass through the field of view: shutter budget and downlink capacity are spent on frames scientists cannot use.
We ask whether a learned attitude policy can outperform a deterministic baseline by using simulated vision feedback to time captures and point for image quality under the same mission geometry and environment setup.
A cloud-aware 2D simulation couples attitude dynamics, dual-camera ray tracing, and safety arbitration with maximum \emph{a posteriori} policy optimization (MPO) over reaction-wheel torques; evaluation uses matched baseline-overflight episodes (notebook~07) as the reference controller.
Qualitative multi-seed rollouts indicate that the learned policy prioritizes novel, high-value targets and avoids blind shuttering along the target stripe, whereas the baseline images every target vector without discriminating sensor content---quantitative yield comparisons are reported in the Results section.
The study is limited to a single-pass, simplified 2D orbital model; attitude safety constraints are enforced, but flight hardware, multi-orbit planning, and full three-dimensional operations are out of scope.
